\subsection{商群}

命题 16 设 G 是一个拓扑群，H 是 G 的一个正规子群。则 G 的拓扑经 H 的商与 G/H 的群结构相容。

若 \(x \to \dot{x}\) 是 G 到 G/H 上的典范映射，则我们要证明的是：\((\dot{x}, \dot{y}) \to \dot{x}\dot{y}^{-1}\) 是 (G/H) \(\times\) (G/H) 到 G/H 中的一个连续映射。由于等价关系 \(x^{-1}y \in H\) 是开的（第 4 目引理 2），只需证明 \((x,y) \to \dot{x}\dot{y}^{-1}\) 是 G \(\times\) G 到 G/H 中的一个连续映射（第一章 § 5 第 3 目命题 8 的推论，以及 § 3 第 4 目命题 6）。但 \((x,y) \to \dot{x}\dot{y}^{-1}\) 是连续映射 \(x \to \dot{x}\) 与 \((x,y) \to xy^{-1}\) 的复合，因而是连续的。

在后续内容中，每当我们把拓扑群 G 的商群 G/H 作为拓扑群考虑时，除相反的情况有明确说明外，总应理解为 G/H 的拓扑是 G 的拓扑经 H 的商。

命题 17 设 \(\varphi\) 是拓扑群 \(\mathbf{G}\) 到商群 \(\mathrm{G / H}\) 上的典范映射。若 \(\mathfrak{B}\) 是 \(e\) 在 \(\mathbf{G}\) 中的一个基本邻域系，则 \(\varphi(\mathfrak{B})\) 是单位元 \(\varphi(e)\)（\(\mathrm{G / H}\) 的）的一个基本邻域系。

这是第一章 § 5 第 3 目命题 5 的一个特殊情形。

命题 13 与 14 对商群特别地给出：

命题 18 设 G 是一个拓扑群，H 是 G 的一个正规子群。

\begin{enumerate}
\def\labelenumi{\alph{enumi})}
\item
  商群 G/H 是 Hausdorff 的，当且仅当 H 在 G 中是闭的。
\item
  商群 G/H 是离散的，当且仅当 H 在 G 中是开的。
\end{enumerate}

若 G 是一个拓扑群，N 是 \(\{e\}\) 在 G 中的闭包，则 N 是 G 的一个闭正规子群（第 1 目命题 1），因而 G/N 是 Hausdorff 的；G/N 称为 G 的伴随 Hausdorff 群。

命题 19 若 H 是拓扑群 G 的一个离散正规子群，则 G/H 与 G 局部同构。

设 \(\mathbf{V}\) 是 \(e\) 在 \(\mathbf{G}\) 中的一个邻域，它不含 \(\mathbf{H}\) 中除 \(e\) 外的任何点，又设 \(\mathbf{W}\) 是 \(e\) 在 \(\mathbf{G}\) 中的一个对称开邻域，使得 \(\mathbf{W}^2 \subset \mathbf{V}\)。则 \(\mathbf{W}\) 上典范映射 \(\varphi\)（\(\mathbf{G}\) 到 \(\mathbf{G} / \mathbf{H}\) 上的）的限制是单射；因为若 \(x, y \in \mathbf{W}\) 且 \(\varphi(x) = \varphi(y)\)，则 \(x^{-1}y \in \mathbf{W}^2 \subset \mathbf{V}\) 且 \(x^{-1}y \in \mathbf{H}\)，于是 \(x = y\)。由命题 17 可知 \(\varphi\) 在 \(\mathbf{W}\) 上的限制是 \(\mathbf{W}\) 到 \(\varphi(\mathbf{W})\) 上的一个同胚；又因 \(\varphi(xy) = \varphi(x)\varphi(y)\) 对所有 \(x, y \in \mathbf{W}\) 成立，我们断定 \(\mathbf{G}\) 与 \(\mathbf{G} / \mathbf{H}\) 局部同构（§ 1 第 3 目命题 3）。

